\subsection{EXTERIOR ALGEBRA OF A DIRECT SUM. EXTERIOR ALGEBRA OF A GRADED MODULE}

Let A be a commutative ring, \(M = \bigoplus_{\lambda \in L} M_{\lambda}\) the direct sum of a family of a family of A-modules and \(j_{\lambda}: M_{\lambda} \to M\) the canonical injection; an A-homomorphism of graded algebras \(\bigwedge(j_{\lambda}): \bigwedge(M_{\lambda}) \to \bigwedge(M)\) is derived. Since \(\bigwedge(M)\) is anticommutative, Proposition 10 of §4, no. 7 (or if need be generalized to the case where L is infinite, cf.~§4, no. 8, Remark 1) can be applied to the homomorphisms \(\bigwedge(j_{\lambda})\) ; then there exists a unique algebra homomorphism

\[
g \colon \bigotimes_ {\lambda \in L} ^ {g} \Lambda (M _ {\lambda}) \to \Lambda (M)\tag{14}
\]

(also denoted by \(g_{\mathbf{M}}\) ) such that \(\bigwedge (j_{\lambda}) = g \circ f_{\lambda}\) , where

\[
f _ {\lambda}: \bigwedge (\mathrm{M} _ {\lambda}) \rightarrow \bigotimes_ {\lambda \in L} ^ {g} \bigwedge (\mathrm{M} _ {\lambda})
\]

denotes the canonical homomorphism.

PROPOSITION 10. The canonical homomorphism \(g\) (formula (14)) is a graded algebra isomorphism.

To prove that \(g\) is bijective, we define a graded algebra homomorphism

\[
h: \bigwedge (\mathbf {M}) \rightarrow \bigotimes_ {\lambda \in L} ^ {\mathfrak {g}} \bigwedge (\mathbf {M} _ {\lambda})\tag{15}
\]

such that \(g \circ h\) and \(h \circ g\) are respectively the identity mappings on \(\bigwedge(\mathbf{M})\) and \(\bigotimes_{\lambda \in L} \bigwedge(\mathbf{M}_{\lambda})\) . For each \(\lambda \in L\) , consider the composite linear mapping

\[
u _ {\lambda}: \mathrm{M} _ {\lambda} \xrightarrow {\phi_ {\mathrm{M} _ {\lambda}} ^ {n}} \bigwedge (\mathrm{M} _ {\lambda}) \xrightarrow {f _ {\lambda}} \mathbf {g} _ {\lambda \in L} \bigotimes \bigwedge (\mathrm{M} _ {\lambda}).
\]

There exists one and only one A-linear mapping \(u: \mathbf{M} \to {}^{\mathfrak{g}} \bigotimes_{\lambda \in L} \bigwedge (\mathbf{M}_{\lambda})\) such that \(u \circ j_{\lambda} = u_{\lambda}\) for all \(\lambda \in L\) . The skew tensor product \({}^{\mathfrak{g}} \bigotimes_{\lambda \in L} \bigwedge (\mathbf{M}_{\lambda})\) is an alternating algebra: for, for L finite, this follows from § 4, no. 9, Proposition 14 and, for L arbitrary, this follows from the definition of this product given in § 4, no. 8, Remark 1 and the fact that a direct limit of alternating graded algebras is alternating. As also \(u(\mathbf{M})\) is contained in the submodule of elements of degree 1 in the graded algebra \(\bigotimes_{\lambda \in L} \bigwedge (\mathbf{M}_{\lambda})\) , it follows from no. 1, Proposition 1 and Remark 1 that there exists a unique graded algebra homomorphism.

\[
h: \bigwedge (\mathbf {M}) \rightarrow \bigotimes_ {\lambda \in L} ^ {\mathfrak {g}} \bigwedge (\mathbf {M} _ {\lambda})
\]

such that \(h \circ \phi_{\mathbf{M}}'' = u\) . The verification of the fact that \(g \circ h\) and \(h \circ g\) are the identity mappings is then performed as in §6, no. 6, Proposition 9 replacing \(\mathbf{S}\) by \(\bigwedge\) and \(\phi'\) by \(\phi''\) .

Remark. Let \(N = \bigoplus_{\lambda \in L} N_{\lambda}\) be the direct sum of another family of A-modules with L as indexing set and, for all \(\lambda \in L\) , let \(v_{\lambda}: M_{\lambda} \to N_{\lambda}\) be an A-linear mapping, whence there is an A-linear mapping \(v = \bigoplus_{\lambda} v_{\lambda}: M \to N\) (II, §1, no. 6, Proposition 7). Then the diagram

\[
\begin{array}{c c c} \mathbf {g} _ {\lambda \in L} ^ {\otimes} \wedge (\mathrm{M} _ {\lambda}) & \xrightarrow {g _ {\mathrm{M}}} & \wedge (\mathrm{M}) \\ \bigotimes_ {\lambda \in L} \wedge (v _ {\lambda}) \Big \downarrow & & \Big \downarrow \wedge (v) \\ \mathbf {g} _ {\lambda \in L} ^ {\otimes} \wedge (\mathrm{N} _ {\lambda}) & \xrightarrow {g _ {\mathrm{N}}} & \wedge (\mathrm{N}) \end{array}
\]

is commutative (cf.~§ 4, no. 5, Corollary to Proposition 8).

The sub-A-module of \(\mathbf{g}_{\lambda \in L}^{\otimes} \bigwedge (\mathbf{M}_{\lambda})\) with which \(\bigwedge^{n}(\mathbf{M})\) is identified by means of the isomorphism \(g\) can be described more precisely. For every finite subset \(J\) of \(L\) , we write \(\mathrm{E}_{J} = \mathbf{g}_{\lambda \in J}^{\otimes} \bigwedge (\mathbf{M}_{\lambda})\) , so that \(\mathbf{g}_{\lambda \in L}^{\otimes} \bigwedge (\mathbf{M}_{\lambda}) = \lim_{\lambda \to 0} \mathrm{E}_{J}\) relative to the directed set \(\mathfrak{F}(L)\) of finite subsets of \(L\) , by definition (§ 4, no. 8, Remark 1). For every family \(\nu = (n_{\lambda}) \in \mathbf{N}^{(L)}\) (therefore with finite support) such that \(\sum_{\lambda \in L} n_{\lambda} = n\) and every finite subset \(J\) of \(L\) containing the support of the family \(\nu\) , we write

\[
\bigwedge^ {J, v} (\mathbf {M}) = \bigotimes_ {\lambda \in J} \bigwedge^ {n _ {\lambda}} (\mathbf {M} _ {\lambda})\tag{16}
\]

so that the submodule \(\mathbf{E}_{\mathbf{J},n}\) of elements of degree \(n\) in \(\mathbf{E}_{\mathbf{J}}\) is the direct sum of the \(\bigwedge^{\mathbf{J},\nu}(\mathbf{M})\) for all families \(\nu\) of support contained in \(\mathbf{J}\) and such that

\[
\sum_ {\lambda \in L} n _ {\lambda} = n
\]

(§ 4, no. 7, Proposition 10 and no. 8). By way of convention we write \(\bigwedge^{J,\nu}(\mathbf{M}) = \{0\}\) for the families \(\nu\) whose support is not contained in \(J\) ; then

\(E_{J,n}\) can also be called the direct sum of all the \(\bigwedge^{J,\nu}(M)\) , where \(\nu\) runs through the set \(H_n\) of all families \(\nu = (n_\lambda)_{\lambda \in L}\) such that \(\sum_{\lambda \in L} n_\lambda = n\) . Since \(\bigwedge^0(M_\lambda)\) is identified with A, clearly also for two finite subsets \(J \subset J'\) of L and a family \(\nu\) of support contained in J, the canonical mapping \(\bigwedge^{J,\nu}(M) \to \bigwedge^{J',\nu}(M)\) (restriction of the canonical mapping \(E_J \to E_{J'}\) to \(\bigwedge^{J,\nu}(M)\) ) is bijective. If we write, for all \(\nu \in H_n\) ,

\[
\bigwedge^ {\nu} (\mathbf {M}) = \varinjlim \bigwedge^ {J, \nu} (\mathbf {M})\tag{17}
\]

it is then seen that, taking account of II, § 6, no. 2, Proposition 5, we have:

COROLLARY. The A-module \(\bigwedge^{n}(\mathbf{M})\) is the image under isomorphism (14) of the submodule of \(\bigotimes_{\lambda \in L} \bigwedge (\mathbf{M}_{\lambda})\) the direct sum of the submodules \(\bigwedge^{\nu}(\mathbf{M})\) for all families \(\nu = (n_{\lambda}) \in \mathbf{N}^{(L)}\) such that \(\sum_{\lambda \in L} n_{\lambda} = n\) ; if \(J\) is the support of \(\nu, \bigwedge^{\nu}(\mathbf{M})\) is canonically isomorphic to \(\bigotimes_{\lambda \in J} \bigwedge^{n_{\lambda}}(\mathbf{M}_{\lambda})\) .

In general \(\bigwedge^{\nu}(\mathbf{M}),\bigotimes_{\lambda \in J}\bigwedge^{n_{\lambda}}(\mathbf{M}_{\lambda})\) and their image in \(\bigwedge^n (\mathbf{M})\) are identified. With this convention:

PROPOSITION 11. Let \(\Delta\) be a commutative monoid, \(\mathbf{M}\) a graded A-module of type \(\Delta\) and \((\mathbf{M}_{\alpha})_{\alpha \in \Delta}\) its graduation. For every ordered pair \((\alpha, n) \in \Delta \times \mathbf{N}\) , let \(\bigwedge^{\alpha, n}(\mathbf{M})\) be the submodule of \(\bigwedge^n (\mathbf{M})\) the direct sum of the submodules \(\bigotimes_{\lambda \in J} \bigwedge^{n_{\lambda}}(\mathbf{M}_{\alpha_{\lambda}})\) , where \((n_{\lambda})_{\lambda \in L}\) runs through the set of families of integers \(\geqslant 0\) such that \(\sum_{\lambda \in L} n_{\lambda} = n\) , \(J\) is its support and, for each \((n_{\lambda})\) , \((\alpha_{\lambda})_{\lambda \in J}\) runs through the set of families of \(\Delta^J\) such that \(\sum_{\lambda \in J} \alpha_{\lambda} = \alpha\) . Then \((\bigwedge^{\alpha, n}(\mathbf{M}))_{(\alpha, n) \in \Delta \times N}\) is the only graduation of type \(\Delta \times N\) compatible with the algebra structure on \(\bigwedge(\mathbf{M})\) and which induces on \(\mathbf{M} = \bigwedge^1 (\mathbf{M})\) the given graduation.

The fact that \(\bigwedge (\mathbf{M})\) is the direct sum of the \(\bigwedge^{\alpha ,n}(\mathbf{M})\) follows from the Corollary to Proposition 10; the rest of the proof is identical with the end of the proof of §5, no. 5, Proposition 7.

